\section{Introduction}\label{sec:introduction}

Long-term, large-scale vineyard monitoring is essential for understanding disease dynamics, identifying spatial and temporal patterns, and supporting sustainable disease management strategies. Among the diseases affecting grapevine production worldwide, esca and other grapevine trunk diseases (GTDs) remain among the major challenges facing viticulture. These diseases cause substantial yield losses, vine decline and mortality, leading to significant economic consequences \citep{bertsch_grapevine_2013,gramaje_managing_2018,hofstetter_what_2012}. Across Europe, esca has been increasingly reported over recent decades. \citep{bruez_overview_2012,dewasme_long-term_2022,etienne_exploring_2024} A similar increase has been reported globally \citep{kenfaoui_panoramic_2022}.

Esca is a complex disease syndrome characterized by necrotic lesions in woody tissues and foliar symptoms that can culminate in complete canopy loss \citep{fischer_grapevine_2019,ouadi_ecophysiological_2019}. Despite extensive research, the mechanisms underlying symptom expression remain incompletely understood. In particular, the precise role of pathogenic microorganisms in foliar symptom development is still debated \citep{bruez_major_2020,del_frari_characterization_2019,monod_landscape-scale_2025}. Current evidence suggests that symptom expression results from interactions among multiple biotic and abiotic factors acting as predisposing conditions or triggering events. Host-related factors reported to influence disease expression include cultivar susceptibility \citep{chacon-vozmediano_cultivar_2021,gastou_large_2024,guerin-dubrana_grapevine_2019}, vineyard age \citep{etienne_exploring_2024,fussler_characterization_2008}, and nitrogen fertilization practices \citep{dellacqua_nitrogen_2025}. Environmental conditions play an equally prominent role. Esca incidence has been linked to temperature regimes, precipitation patterns, drought events, and other climatic variables affecting plant water status and physiological functioning \citep{bortolami_grapevines_2021,calzarano_rainfall_2018,serra_expression_2018,surico_epidemiology_2000}. Likewise, soil characteristics influencing water availability have been associated with increased symptom expression \citep{monod_identifying_2025}. Collectively, these findings highlight the importance of plant-environment interactions in shaping annual esca incidence.

A major challenge in studying esca epidemiology is the strong year-to-year variability in symptom expression. Esca incidence can fluctuate substantially between consecutive seasons \citep{calzarano_rainfall_2018,etienne_exploring_2024,etienne_identification_2026}, making short-term observations potentially misleading when assessing long-term disease dynamics. Consequently, long-term monitoring is necessary to distinguish transient fluctuations from genuine epidemiological trends. Furthermore, although climatic factors are frequently associated with symptom expression, their ability to predict disease incidence across sites and years remains insufficiently explored.

To address these questions, we conducted a 21-year survey of a network of vineyard plots representing ten grapevine cultivars across six viticultural regions of Western Switzerland . Esca foliar symptom incidence was monitored annually or biennially throughout the study. This long-term dataset enables the investigation of temporal and spatial patterns of esca expression and the assessment of the predictive value of meteorological variables at the plot level. We selected climatic indicators known to influence grapevine physiology and previously associated with esca incidence, including temperature, atmospheric pressure, relative humidity, radiation, precipitation, evapotranspiration, and vapor pressure deficit. Using a predictive modelling approach applied to climatic time-series data, we tested the extent to which these environmental variables could predict esca incidence at both plot and regional scales. We hypothesized that climatic variables affecting grapevine water balance would explain a significant proportion of interannual variation in esca incidence and could therefore help predict disease occurrence across vineyards and regions.
